\subsection{The exact reduction from first passage to the main theorem}
\label{subsec:Tao-main-reduction}

The source reduction from Proposition~1.11 to Theorems~3.1, 1.6, and~1.3
is v7 lines~535--593, with the definitions and statements at lines~159--206
\cite{Tao2026v7}.  The corresponding v5 loci are lines~529--587 and
159--206; after deletion of whitespace, the two reduction blocks are
identical.  The journal manifestation reproduces the same argument on
physical pages~15--17, with its logarithmic-density and Syracuse setup on
physical pages~3--4 and Proposition~1.11 on physical page~8.  The comparison
below is therefore version-pinned but not a claim that the complete
manifestations are text-identical.

\begin{sourceissue}[defects and suppressed arrows in the printed reduction]
\label{sourceissue:Tao-main-reduction}
Eight points must be separated.

First, v7 lines~591--593 define
\(\widetilde f(x)=\inf_{N\geq x}f(N)\) and then use that lower envelope on
the range \(N\leq x\), where it supplies no comparison.  The displayed
bound is false: on the odd positive integers, take \(f(1)=0\) and
\(f(N)=e^N\) for \(N\geq3\).  The term \(N=1\) contributes one to the
printed left side, while its asserted right side is
\(O(\log x/x^c)\).  Second, the complement of
\(\operatorname{Syr}_{\min}(N)<f(N)\) is
\(\operatorname{Syr}_{\min}(N)\geq f(N)\), not the printed strict
inequality in the other direction; equality cannot be discarded.  Third,
the statement permits real-valued \(f\), while line~591 silently restricts
its codomain to \([0,+\infty)\).

Fourth, line~206 assigns logarithmic density \(2^{-a}\) to the stratum
\(\nu_2(N)=a\) and \(1-2^{-a_0}\) to the union
\(0\leq a\leq a_0\).  The exact values are respectively
\(2^{-(a+1)}\) and \(1-2^{-(a_0+1)}\); the function on the odd coordinate
must moreover be \(M\mapsto f(2^aM)\), not \(f(M)\).

Fifth, the iteration at lines~569--580 omits the possible value \(J=0\).
Sixth, its probabilistic blocks need not be defined at small scales because
the narrow interval can contain no odd integer.  Seventh, the parenthetical
at line~552 says that finite passage is redundant when \(x\geq N_0\); this
is false under the source convention \(\operatorname{Pass}_x=1\) when
\(T_x=+\infty\), since \(1\in E_{N_0}\).  The displayed proof does retain
the finite-passage condition.  Eighth, the substitution at lines~569--574
produces
\(O((\alpha^{j-2}\log y)^{-c})\), not the printed
\(O((\alpha^j\log y)^{-c})\).  Their ratio is the fixed factor
\(\alpha^{2c}\), so the iteration survives only after that factor is
explicitly absorbed.

The same passage also suppresses the indexed first-passage morphism, the
actual multiplicative block cover, its harmonic normalisation, and the
dyadic weight transport.  These are proved below.  None of the repairs is
silently attributed to the source.
\end{sourceissue}

For a real cutoff \(X\geq1\), set
\[
 H_{\rm odd}(X)=
 \sum_{\substack{1\leq M\leq X\\M\ {\rm odd}}}\frac1M,
 \qquad
 H(X)=\sum_{1\leq N\leq X}\frac1N.
\]
The integral test, with the integer cutoffs retained, gives
\begin{equation}
\label{eq:Tao-main-harmonic-asymptotics}
 H_{\rm odd}(X)=\frac12\log X+O(1),
 \qquad H(X)=\log X+O(1).
\end{equation}
For a real threshold \(N_0\geq2\), write
\begin{equation}
\label{eq:Tao-main-EN0}
 E_{N_0}=\{M\in2\mathbb N+1:
        \operatorname{Syr}_{\min}(M)\leq N_0\}.
\end{equation}

\begin{lemma}[the indexed nested-threshold morphism]
\label{lem:Tao-nested-passage-tail}
Let \(1\leq u\leq v\), and put
\[
 \mathcal D_u=\{N\in2\mathbb N+1:T_u(N)<+\infty\}.
\]
For \(N\in\mathcal D_u\), define
\[
 d_{u,v}(N)=T_u(N)-T_v(N)\in\mathbb N,
 \qquad \sigma_d(k)=k+d.
\]
Then
\begin{equation}
\label{eq:Tao-nested-pass-point}
 \operatorname{Pass}_u(N)
 =\operatorname{Syr}^{d_{u,v}(N)}
    (\operatorname{Pass}_v(N)).
\end{equation}
More precisely, for the indexed tails
\[
 \mathcal O_w^N(k)=
 \operatorname{Syr}^k(\operatorname{Pass}_w(N)),
 \qquad k\in\mathbb N,
\]
one has the exact identity
\begin{equation}
\label{eq:Tao-nested-tail-shift}
 \mathcal O_u^N=\mathcal O_v^N\circ\sigma_{d_{u,v}(N)}.
\end{equation}
The shift \(\sigma_d\) is injective, has image \(d+\mathbb N\), has
singleton fibres over that image, and has diagonal equality kernel.  Hence
the value-set of the \(u\)-tail is contained in the value-set of the
\(v\)-tail.  In particular,
\begin{equation}
\label{eq:Tao-nested-EN0-transport}
 T_u(N)<+\infty,\quad \operatorname{Pass}_u(N)\in E_{N_0}
 \quad\Longrightarrow\quad
 \operatorname{Pass}_v(N)\in E_{N_0}.
\end{equation}
\end{lemma}

\begin{proof}
Every iterate at most \(u\) is also at most \(v\), so
\(T_v(N)\leq T_u(N)<+\infty\).  The semigroup law for the iterates gives
\[
 \operatorname{Syr}^{T_u(N)}(N)
 =\operatorname{Syr}^{T_u(N)-T_v(N)}
   (\operatorname{Syr}^{T_v(N)}(N)),
\]
which is \eqref{eq:Tao-nested-pass-point}.  Applying a further \(k\)
iterates gives \eqref{eq:Tao-nested-tail-shift}.  If the later tail from
\(\operatorname{Pass}_u(N)\) reaches \(N_0\), the containing tail from
\(\operatorname{Pass}_v(N)\) reaches the same value, proving
\eqref{eq:Tao-nested-EN0-transport}.  No injectivity of
\(\mathcal O_v^N\) is used or asserted: a periodic orbit can repeat values
at different indices.
\end{proof}

For every scale \(s\) for which the odd logarithmic block is nonempty,
retain
\[
 \mathcal R_s=(2\mathbb N+1)\cap[s,s^\alpha],
 \qquad H_s=\sum_{N\in\mathcal R_s}\frac1N,
 \qquad
 \mathbb P(\mathbf N_s=N)=\frac1{H_sN}.
\]
Define the numerical event probability
\begin{equation}
\label{eq:Tao-main-bs}
 b_{N_0}(s)=\mathbb P\left(
  T_s(\mathbf N_{s^\alpha})<+\infty,
  \ \operatorname{Pass}_s(\mathbf N_{s^\alpha})\in E_{N_0}
 \right).
\end{equation}

\begin{proposition}[the exact block recurrence]
\label{prop:Tao-main-block-recurrence}
There are absolute constants \(X_*,C,c>0\) such that, whenever
\(s\geq X_*\),
\begin{equation}
\label{eq:Tao-main-block-recurrence}
 b_{N_0}(s^\alpha)
 \geq b_{N_0}(s)-C(\log s)^{-c}.
\end{equation}
The cutoff \(X_*\) is chosen so that Proposition~1.11 applies and every
block occurring here contains an odd integer.
\end{proposition}

\begin{proof}
On the probability space supporting \(\mathbf N_{s^{\alpha^2}}\), let
\[
 A_s=\{T_s(\mathbf N_{s^{\alpha^2}})<+\infty,
 \ \operatorname{Pass}_s(\mathbf N_{s^{\alpha^2}})\in E_{N_0}\}.
\]
Lemma~\ref{lem:Tao-nested-passage-tail}, with \(u=s\) and
\(v=s^\alpha\), gives
\[
 A_s\subseteq
 \{T_{s^\alpha}(\mathbf N_{s^{\alpha^2}})<+\infty,
 \ \operatorname{Pass}_{s^\alpha}
       (\mathbf N_{s^{\alpha^2}})\in E_{N_0}\}.
\]
The probability of the event on the right is
\(b_{N_0}(s^\alpha)\).  Proposition~1.11 and the event inequality for its
unhalved total-variation convention therefore give
\begin{align*}
 b_{N_0}(s^\alpha)
 &\geq\mathbb P(A_s)\\
 &\geq
 \mathbb P(\operatorname{Pass}_s(\mathbf N_{s^{\alpha^2}})
             \in E_{N_0})-Cs^{-c}\\
 &\geq
 \mathbb P(\operatorname{Pass}_s(\mathbf N_{s^\alpha})
             \in E_{N_0})
       -C(\log s)^{-c}-Cs^{-c}\\
 &\geq b_{N_0}(s)-C'(\log s)^{-c}.
\end{align*}
Increasing \(C\) proves \eqref{eq:Tao-main-block-recurrence}.  This uses
the two marginal passage-location laws and failure control; it does not
construct a coupling.
\end{proof}

\begin{theorem}[repaired exact form of Tao's Theorem~3.1]
\label{thm:Tao-main-alt-repaired}
There are absolute constants \(C,c>0\) such that, for all real
\(N_0\geq2\) and all real \(X\geq2\),
\begin{equation}
\label{eq:Tao-main-alt-odd}
 \sum_{\substack{M\in2\mathbb N+1,\ M\leq X\\
       \operatorname{Syr}_{\min}(M)>N_0}}
       \frac1M
 \leq C\frac{\log X}{(\log N_0)^c},
\end{equation}
and
\begin{equation}
\label{eq:Tao-main-alt-collatz}
 \sum_{\substack{N\in\mathbb N+1,\ N\leq X\\
       \operatorname{Col}_{\min}(N)>N_0}}
       \frac1N
 \leq C\frac{\log X}{(\log N_0)^c}.
\end{equation}
Consequently the corresponding good sets have, uniformly in \(X\), odd
relative and full logarithmic probabilities at least
\(1-O((\log N_0)^{-c})\).
\end{theorem}

\begin{proof}
Choose \(X_*\) as in
Proposition~\ref{prop:Tao-main-block-recurrence}.  We first prove the block
bound
\begin{equation}
\label{eq:Tao-main-block-bad-probability}
 \mathbb P(\operatorname{Syr}_{\min}(\mathbf N_s)>N_0)
 \ll(\log N_0)^{-c}.
\end{equation}
For bounded \(N_0\), this follows after increasing the implied constant.
Thus assume \(N_0\) is large enough that
\(N_0^{1/\alpha^3}\geq X_*\).

If \(s\leq N_0^{1/\alpha}\), every integer in
\([s,s^\alpha]\) is at most \(N_0\), so the probability in
\eqref{eq:Tao-main-block-bad-probability} is exactly zero.  This is the
omitted \(J=0\) case.

Now suppose \(s>N_0^{1/\alpha}\).  Let \(J\geq1\) be least such that
\[
 y=s^{\alpha^{-J}}<N_0^{1/\alpha}.
\]
Minimality gives
\begin{equation}
\label{eq:Tao-main-y-bounds}
 y^\alpha\geq N_0^{1/\alpha},
 \qquad y\geq N_0^{1/\alpha^2},
 \qquad y^\alpha<N_0.
\end{equation}
Put \(r_0=y^{1/\alpha}\).  Then
\(r_0\geq N_0^{1/\alpha^3}\geq X_*\), and the input in
\(b_{N_0}(r_0)\) is \(\mathbf N_y\).  On finite passage below \(r_0\),
the passage location is less than \(N_0\), so Proposition~1.11 gives
\begin{equation}
\label{eq:Tao-main-base-event}
 b_{N_0}(r_0)\geq1-Cr_0^{-c}.
\end{equation}
For \(j=1,\ldots,J\), apply
\eqref{eq:Tao-main-block-recurrence} at the exact scale
\(y^{\alpha^{j-2}}\).  The accumulated logarithmic error is
\begin{align}
 \sum_{j=1}^J
  (\alpha^{j-2}\log y)^{-c}
 &\leq (\log y)^{-c}
     \sum_{j=1}^{\infty}\alpha^{-c(j-2)}\notag\\
 &=\frac{\alpha^c}{1-\alpha^{-c}}(\log y)^{-c}
 \ll(\log N_0)^{-c}.
\label{eq:Tao-main-geometric-error}
\end{align}
Here \eqref{eq:Tao-main-y-bounds} was used only after the exact factor in
the first line was retained.  Equations
\eqref{eq:Tao-main-base-event} and
\eqref{eq:Tao-main-geometric-error} yield
\[
 b_{N_0}(s^{1/\alpha})\geq1-C(\log N_0)^{-c}.
\]
Its input variable is exactly \(\mathbf N_s\).  Whenever its event holds,
the orbit of \(\mathbf N_s\) reaches \(E_{N_0}\), proving
\eqref{eq:Tao-main-block-bad-probability}.

For \(s\geq2\), the integral test gives \(H_s\ll\log s\), with an
absolute constant because \(\alpha=1001/1000\) is fixed.  An applied cover
scale with \(1<s<2\) has no odd integer in \([s,s^\alpha]\), so its local
sum is zero; every nonempty cover block to which the estimate is applied
satisfies \(s\geq2\).  Hence
\begin{equation}
\label{eq:Tao-main-local-harmonic-bound}
 \sum_{\substack{M\in(2\mathbb N+1)\cap[s,s^\alpha]\\
          \operatorname{Syr}_{\min}(M)>N_0}}
          \frac1M
 \ll\frac{\log s}{(\log N_0)^c}.
\end{equation}
For a target \(X>N_0\), set
\(z_k=X^{\alpha^{-k}}\), and let \(L\geq1\) be least such that
\(z_L\leq N_0\).  Then
\[
 z_{L-1}>N_0,
 \qquad z_L>N_0^{1/\alpha},
 \qquad
 [z_L,X]=\bigcup_{k=1}^L[z_k,z_k^\alpha].
\]
Every bad starting value is greater than \(N_0\), since its orbit contains
the starting value.  Thus the leftover interval \([1,z_L]\) is bad-free,
and endpoint duplication in the displayed cover can only enlarge an upper
bound.  Finally,
\begin{equation}
\label{eq:Tao-main-cover-log-sum}
 \sum_{k=1}^L\log z_k
 =\log X\sum_{k=1}^L\alpha^{-k}
 \leq\frac{\log X}{\alpha-1}.
\end{equation}
Summing \eqref{eq:Tao-main-local-harmonic-bound} proves
\eqref{eq:Tao-main-alt-odd}.  If \(X\leq N_0\), its left side is zero.

For the all-integer assertion, write every \(N\) uniquely as
\(N=2^aM\), with \(a\in\mathbb N\) and \(M\) odd.  The exact minimum
identity proved again below gives
\begin{equation}
\label{eq:Tao-main-fixed-dyadic-sum}
 \sum_{\substack{N\leq X\\
       \operatorname{Col}_{\min}(N)>N_0}}\frac1N
 =\sum_{a\geq0}2^{-a}
  \sum_{\substack{M\leq X/2^a,\ M\ {\rm odd}\\
       \operatorname{Syr}_{\min}(M)>N_0}}\frac1M.
\end{equation}
When \(X/2^a<2\), the inner bad sum vanishes because \(N_0\geq2\).
For the remaining terms, \eqref{eq:Tao-main-alt-odd} and the truncated,
termwise nonnegative estimate
\[
 \sum_{\substack{a\geq0\\X/2^a\geq2}}
     2^{-a}\log(X/2^a)
 \leq \log X\sum_{a\geq0}2^{-a}
 \leq2\log X
\]
prove \eqref{eq:Tao-main-alt-collatz}.  The probability statements follow
from \eqref{eq:Tao-main-harmonic-asymptotics}; no identity of unnormalised
and normalised harmonic mass is being asserted.
\end{proof}

\begin{theorem}[repaired exact form of Tao's Theorem~1.6]
\label{thm:Tao-main-syr-repaired}
Let \(f:2\mathbb N+1\to\mathbb R\) satisfy
\(f(M)\to+\infty\).  Then
\begin{equation}
\label{eq:Tao-main-syr-strict}
 \operatorname{Syr}_{\min}(M)<f(M)
\end{equation}
for a set of odd positive integers of relative odd logarithmic density one.
Equivalently, the set of odd \(M\) on which
\(\operatorname{Syr}_{\min}(M)\geq f(M)\) has relative odd logarithmic
density zero.
\end{theorem}

\begin{proof}
Fix an integer \(K\geq2\).  Since \(f(M)\to+\infty\), there is a finite
cutoff \(Y_K\) such that
\[
 f(M)>K\qquad(M\geq Y_K,\ M\ {\rm odd}).
\]
Let
\(D_f=\{M\ {\rm odd}:\operatorname{Syr}_{\min}(M)\geq f(M)\}\).
For every \(X\), the exact tail inclusion gives
\begin{align}
 \sum_{\substack{M\leq X\\M\in D_f}}\frac1M
 &\leq H_{\rm odd}(Y_K)
 +\sum_{\substack{M\leq X,\ M\ {\rm odd}\\
          \operatorname{Syr}_{\min}(M)>K}}\frac1M\notag\\
 &\leq H_{\rm odd}(Y_K)
 +C\frac{\log X}{(\log K)^c}.
\label{eq:Tao-main-fixed-K-tail}
\end{align}
Divide first by \(H_{\rm odd}(X)\), let \(X\to+\infty\) with \(K\)
fixed, and only then let \(K\to+\infty\).  Equation
\eqref{eq:Tao-main-harmonic-asymptotics} gives
\[
 \limsup_{X\to\infty}
 \frac{\sum_{M\leq X,\ M\in D_f}1/M}{H_{\rm odd}(X)}
 \leq\frac{2C}{(\log K)^c}\longrightarrow0.
\]
This proves the strict inequality in \eqref{eq:Tao-main-syr-strict}, includes
the equality part of its complement, permits arbitrary nonmonotone real
\(f\), and confines every negative or otherwise exceptional value of \(f\)
to the finite prefix.  No tail lower envelope is used.
\end{proof}

\begin{proposition}[the exact dyadic stratum morphism]
\label{prop:Tao-main-dyadic-morphism}
Let
\[
 \mathcal O=2\mathbb N+1,
 \qquad
 \mathcal S_a=\{N\in\mathbb N+1:\nu_2(N)=a\}.
\]
For each \(a\in\mathbb N\), the maps
\[
 \iota_a:\mathcal O\longrightarrow\mathcal S_a,
 \quad \iota_a(M)=2^aM,
 \qquad
 \pi_a:\mathcal S_a\longrightarrow\mathcal O,
 \quad \pi_a(N)=N/2^a
\]
are two-sided inverses.  Their fibres are singletons, their equality kernels
are diagonal, and they lose no starting-number coordinate information.
Globally,
\begin{equation}
\label{eq:Tao-main-global-dyadic-bijection}
 \mathbb N+1\longleftrightarrow
 \coprod_{a\in\mathbb N}\{a\}\times\mathcal O,
 \qquad
 N\longmapsto(\nu_2(N),N/2^{\nu_2(N)})
\end{equation}
is a bijection with inverse \((a,M)\mapsto2^aM\).

For every \(A\subseteq\mathcal O\) and real \(X\geq1\), the cutoff and
harmonic weight transport exactly as
\begin{equation}
\label{eq:Tao-main-dyadic-weight}
 \sum_{\substack{N\in\iota_a(A)\\N\leq X}}\frac1N
 =2^{-a}\sum_{\substack{M\in A\\M\leq X/2^a}}\frac1M.
\end{equation}
In particular,
\begin{equation}
\label{eq:Tao-main-dyadic-density}
 \sum_{\substack{N\leq X\\\nu_2(N)=a}}\frac1N
 =2^{-a}H_{\rm odd}(X/2^a)
 =2^{-(a+1)}\log X+O_a(1).
\end{equation}
Finally,
\begin{equation}
\label{eq:Tao-main-minimum-identity}
 \operatorname{Col}_{\min}(2^aM)
 =\operatorname{Syr}_{\min}(M).
\end{equation}
\end{proposition}

\begin{proof}
The defining condition \(\nu_2(N)=a\) says uniquely that \(N=2^aM\)
with \(M\) odd, which proves every bijection and fibre assertion.  Replacing
\(N\) by \(2^aM\) in the finite sum proves
\eqref{eq:Tao-main-dyadic-weight}; then
\eqref{eq:Tao-main-harmonic-asymptotics} proves
\eqref{eq:Tao-main-dyadic-density}.

It remains to prove the minimum identity rather than infer it from notation.
The first \(a\) Collatz steps starting at \(2^aM\) are divisions by two and
terminate at \(M\); every value on that initial segment is at least \(M\).
If \(M_j\) is an odd Syracuse state and
\(q_j=\nu_2(3M_j+1)\), then the corresponding Collatz segment is
\[
 M_j,\ 2^{q_j}M_{j+1},\ 2^{q_j-1}M_{j+1},\ldots,
 2M_{j+1},\ M_{j+1},
 \qquad M_{j+1}=\operatorname{Syr}(M_j).
\]
Every suppressed intermediate value is at least its terminal odd value.
Thus the minimum over the complete Collatz orbit is exactly the minimum over
the odd Syracuse states, proving \eqref{eq:Tao-main-minimum-identity}.
The map is not time-preserving.  The full Collatz time coordinate is
reconstructed by expanding each Syracuse step using
\(q_j=\nu_2(3M_j+1)\), which is determined by its odd source state \(M_j\).
\end{proof}

\begin{theorem}[repaired exact form of Tao's Theorem~1.3]
\label{thm:Tao-main-repaired}
Let \(f:\mathbb N+1\to\mathbb R\) satisfy \(f(N)\to+\infty\).  Then
\begin{equation}
\label{eq:Tao-main-collatz-strict}
 \operatorname{Col}_{\min}(N)<f(N)
\end{equation}
for a set of positive integers of logarithmic density one.
\end{theorem}

\begin{proof}
For each fixed \(a\in\mathbb N\), define
\[
 g_a(M)=f(2^aM)\qquad(M\in\mathcal O).
\]
Then \(g_a(M)\to+\infty\), so
Theorem~\ref{thm:Tao-main-syr-repaired} says that
\[
 B_a=\{M\in\mathcal O:
       \operatorname{Syr}_{\min}(M)\geq g_a(M)\}
\]
has harmonic mass \(o(\log T)\) below \(T\).  By
Proposition~\ref{prop:Tao-main-dyadic-morphism}, the bad set for
\eqref{eq:Tao-main-collatz-strict} in \(\mathcal S_a\) is exactly
\(\iota_a(B_a)\), and its mass below \(X\) is
\[
 2^{-a}\sum_{\substack{M\in B_a\\M\leq X/2^a}}\frac1M=o(\log X)
\]
for this fixed \(a\).

Fix \(A\in\mathbb N\).  Splitting the global bad mass into the first
\(A+1\) strata and the remaining tail gives
\begin{align}
 \sum_{\substack{N\leq X\\
   \operatorname{Col}_{\min}(N)\geq f(N)}}\frac1N
 &\leq\sum_{a=0}^A2^{-a}
   \sum_{\substack{M\in B_a\\M\leq X/2^a}}\frac1M
 +\sum_{\substack{N\leq X\\\nu_2(N)>A}}\frac1N,\notag\\
 \sum_{\substack{N\leq X\\\nu_2(N)>A}}\frac1N
 &=2^{-(A+1)}H(X/2^{A+1}).
\label{eq:Tao-main-dyadic-tail-limit}
\end{align}
The second identity follows from the unique formula
\(N=2^{A+1}Q\), with \(Q\) now an arbitrary positive integer.  Divide by
\(H(X)\) and let \(X\to+\infty\) at fixed \(A\).  The finite first sum
vanishes in the limit and
\eqref{eq:Tao-main-harmonic-asymptotics} gives
\[
 \limsup_{X\to\infty}
 \frac{\sum_{N\leq X,\ \operatorname{Col}_{\min}(N)\geq f(N)}1/N}
      {H(X)}
 \leq2^{-(A+1)}.
\]
Only now send \(A\to\infty\).  The bad logarithmic density is zero, which
proves the strict inequality \eqref{eq:Tao-main-collatz-strict}.  No
uniformity in \(a\), and no interchange of the two limits, has been used.
\end{proof}

As an independent cross-check, one may apply the fixed-\(K\) argument in
\eqref{eq:Tao-main-fixed-K-tail} directly to the Collatz estimate
\eqref{eq:Tao-main-alt-collatz}.  This gives
Theorem~\ref{thm:Tao-main-repaired} without passing through the dyadic
strata, and yields the same strict bad-set complement.

\begin{nonclaim}
Theorems~\ref{thm:Tao-main-alt-repaired},
\ref{thm:Tao-main-syr-repaired}, and~\ref{thm:Tao-main-repaired} prove the
exact repaired chain
\[
 \text{Proposition~1.11}\Longrightarrow\text{Theorem~3.1}
 \Longrightarrow\text{Theorem~1.6}
 \Longrightarrow\text{Theorem~1.3}.
\]
They prove logarithmic-density statements, not natural-density statements;
they do not give an absolute bound valid for almost all starting values,
force an individual orbit through a threshold, exclude divergent or
periodic exceptional orbits, prove convergence to one, construct an
 invariant logarithmic measure, or produce a coupling.  The dyadic bijection
preserves all starting-number information---it is a bijection on the
displayed starting coordinates---and preserves the orbit minimum at the
displayed scope.  It is not itself a time-preserving conjugacy; Collatz time
is recovered by the explicit valuation expansion above.
This closes the main-theorem reduction, not the durable whole-paper,
whole-corpus, formalisation, or release gates.
\end{nonclaim}
